\section{Introduction}

Wearable navigation aids promise greater independence for visually impaired
users in unfamiliar indoor spaces, but their sensing stacks are costly. Most
rely on high-resolution RGB or active depth sensors, which continuously capture
identifiable imagery of bystanders and raise privacy concerns when worn indoors.
They also require large hand-labeled datasets to learn which directions are
walkable, and typically assume a fixed, geometrically calibrated camera layout
that breaks as soon as a sensor shifts.

Ultra-low-resolution thermal sensing is an appealing alternative. A
$32\times24$-pixel array costs a few tens of dollars, draws little power, is
invariant to lighting, and is too coarse to resolve identifiable facial
features. These same properties make perception hard: each frame carries little
spatial detail, several cameras are needed for useful coverage, and annotating
thermal data is laborious. We therefore ask whether \emph{directional
navigability}---which directions around the wearer are walkable---can be
estimated from ultra-low-resolution thermal imagery alone, without geometric
calibration and without manual labels.

We answer this affirmatively with a calibration-free, RGB-supervised framework
(Fig.~\ref{fig:pipeline}). During data collection only, a high-resolution RGB
camera rides alongside a wearable rig of three ultra-low-resolution thermal
cameras; an off-the-shelf detector run on the RGB frames automatically generates
per-direction navigability labels, so no manual annotation is needed. At
inference the RGB camera is removed: the model consumes only the three raw
thermal tiles and predicts, for each of the left, middle, and right directions,
whether the space is clear, occupied by a person, or blocked by another object.
Because supervision is transferred through labels rather than pixel-aligned
feature matching, the framework needs no calibration, no pixel-level alignment,
and no fixed correspondence between sensor position and input channel.

Within this framework we study two deployment-critical questions---how to fuse
the three views, and how robust the system is to camera placement---which
motivate our contributions: (i)~a calibration-free framework that transfers
navigability supervision from an RGB detector to a thermal-only model, with no
labels or RGB at deployment; (ii)~a direct comparison of three multi-view fusion
strategies (pixel-, decision-, and feature-level) on ultra-low-resolution
thermal data; and (iii)~an analysis showing that a lightweight
camera-identity correction restores accuracy when camera slots are shuffled,
matching a ground-truth oracle re-sort.

\begin{figure*}[!t]
  \centering
  % Replace the placeholder below with the final pipeline artwork, e.g.:
  % \includegraphics[width=\textwidth]{pipeline_overview.png}
  \figplaceholder[\textwidth]{0.15\textheight}{%
    \textbf{Full-width pipeline figure.}\\[0.3em]
    End-to-end overview: wearable three-camera ultra-low-resolution thermal
    capture (RGB camera used only at data-collection time) $\rightarrow$
    RGB/YOLO automatic per-direction label generation $\rightarrow$ the four
    thermal fusion models (stitched panorama, output fusion, feature fusion,
    feature fusion with camera-identity correction) $\rightarrow$ per-direction
    navigability predictions.}
  \caption{End-to-end pipeline, from multi-view thermal data collection and
  RGB-based automatic labeling to the family of thermal fusion models evaluated
  in this paper. The RGB camera is present only during data collection; at
  inference the models use the three ultra-low-resolution thermal tiles alone.}
  \label{fig:pipeline}
\end{figure*}
